\documentclass[11pt]{article}
\usepackage[utf8]{inputenc}
\usepackage[T1]{fontenc}
\usepackage{amsmath, amssymb, amsthm}
\usepackage{graphicx}
\usepackage{geometry}
\usepackage{hyperref}
\usepackage{listings}
\usepackage{xcolor}
\usepackage{booktabs}
\usepackage{url}

\geometry{margin=1in}

\hypersetup{
  colorlinks=true,
  linkcolor=blue!60!black,
  urlcolor=blue!60!black,
  citecolor=blue!60!black,
}

\lstset{
  basicstyle=\ttfamily\footnotesize,
  frame=single,
  breaklines=true,
  columns=fullflexible,
  backgroundcolor=\color{black!3},
  keywordstyle=\color{blue!60!black}\bfseries,
  commentstyle=\color{green!40!black},
  stringstyle=\color{red!60!black},
}

\newcommand{\vE}{\mathbf{E}}
\newcommand{\vB}{\mathbf{B}}
\newcommand{\vD}{\mathbf{D}}
\newcommand{\vH}{\mathbf{H}}
\newcommand{\vJ}{\mathbf{J}}
\newcommand{\vx}{\mathbf{x}}
\newcommand{\hx}{\hat{x}}
\newcommand{\hy}{\hat{y}}
\newcommand{\hz}{\hat{z}}
\newcommand{\eps}{\varepsilon}

\title{Maxwell\_Penalty: \\
  A 3D Electromagnetic Solver with SBP-SAT Penalty Boundaries}
\author{SMS 751 Lecture Code}
\date{}

\begin{document}

\maketitle

\tableofcontents

\vspace{1em}

\begin{abstract}
This document describes \texttt{Maxwell\_Penalty}, a three-dimensional
finite-difference solver for Maxwell's equations in a linear dispersive
dielectric, coupled to a hyperbolic constraint-damping system. The spatial
discretization uses fourth-order summation-by-parts (SBP) operators with
one-sided closures at every physical face. Boundary conditions are
imposed weakly via simultaneous approximation terms (SAT) — the
penalty method — on all six Cartesian faces of the computational
domain. The time integrator is classical fourth-order Runge--Kutta.
The code includes two analytic sources — a plane wave with a compactly
supported envelope, and a paraxial linearly-polarised Gaussian beam —
and supports an inhomogeneous-permittivity region that acts as a lens
through which the beam propagates. This document covers the physics,
the numerical method, the penalty-boundary treatment, the source models
and their limitations, and how to build, configure, and run the code.
\end{abstract}

\section{Introduction: the physics}
\label{sec:physics}

The code evolves the electromagnetic field $(\vE, \vB)$ in a linear,
possibly inhomogeneous, dielectric. Let $\eps(\vx)$ be the permittivity,
$\mu(\vx)$ the permeability, and $\sigma(\vx)$ the conductivity
(Ohmic damping). Define $\vD = \eps \vE$ and $\vH = \vB / \mu$. In
Gaussian units, Maxwell's equations are
\begin{align}
  \partial_t \vD &= \nabla \times \vH - 4\pi \vJ, \label{eq:ampere}\\
  \partial_t \vB &= -\nabla \times \vE, \label{eq:faraday}\\
  \nabla \cdot \vD &= 4\pi \rho, \label{eq:gauss-e}\\
  \nabla \cdot \vB &= 0, \label{eq:gauss-b}
\end{align}
with Ohm's law $\vJ = \sigma \vE$ for the current density. Equations
\eqref{eq:ampere}--\eqref{eq:faraday} are the time-evolution (\emph{hyperbolic})
equations. Equations \eqref{eq:gauss-e}--\eqref{eq:gauss-b} are the
divergence constraints. In the continuum these constraints are
preserved by the time evolution: if $\vD, \vB$ initially satisfy them,
they satisfy them for all time.

Numerically this is not automatic. Round-off error and truncation error
in the discretized curl operators can seed small constraint violations
that, in a naive scheme, grow without bound. The standard remedy,
following Munz, Omnes, Schneider, and others, is to \emph{extend the
system} with two auxiliary scalar fields $\Psi_D$ and $\Psi_B$ that
act as Lagrange multipliers for the constraints. The extended system
(in vacuum form, generalized below) is
\begin{align}
  \partial_t \vD - c^2 \nabla \Psi_D &= \nabla \times \vH - 4\pi \vJ, \\
  \partial_t \vB - c^2 \nabla \Psi_B &= -\nabla \times \vE, \\
  \partial_t \Psi_D &= \nabla \cdot \vD - 4\pi \rho - \kappa_D \Psi_D, \\
  \partial_t \Psi_B &= \nabla \cdot \vB - \kappa_B \Psi_B.
\end{align}
At $\Psi_D = \Psi_B = 0$ the original Maxwell system is recovered.
For $\kappa_{D}, \kappa_{B} > 0$ the $\Psi$ fields satisfy damped wave
equations whose source is the constraint violation — any residual violation
radiates outward at speed $c$ and is exponentially damped on timescale
$1/\kappa$. This absorbs round-off and truncation noise, and absorbs
the $O(1/(k w_0)^2)$ constraint residual of the paraxial beam source
(Sec.~\ref{sec:gaussian-beam}).

The code evolves nine fields: $(D_x, D_y, D_z, B_x, B_y, B_z, \Psi_D, \Psi_B, \rho)$.
The charge density $\rho$ is carried as an evolved variable for generality,
driven by $\partial_t \rho + \nabla \cdot \vJ = 0$; for the simulations
shipped with the code, $\rho$ and $\vJ$ are initially zero and remain
zero up to round-off.

The material parameters $\eps$, $\mu$, $\sigma$ are stored as auxiliary
gridfunctions in the local inverse form $\mathtt{ieps} = 1/\eps$,
$\mathtt{imu} = 1/\mu$, $\mathtt{sigma} = \sigma$. They are spatially
varying in general; the default configuration places an
elliptical-contour inhomogeneity in $\eps$ to model a lens.

\subsection{The flux form}

The spatial derivatives in the extended system can be written as a
quasi-linear hyperbolic system
\begin{equation}
  \partial_t U + A_x \partial_x U + A_y \partial_y U + A_z \partial_z U
  = S(U, \vx),
  \label{eq:flux-form}
\end{equation}
where $U = (D_x, D_y, D_z, B_x, B_y, B_z, \Psi_D, \Psi_B, \rho)^\top$ and
$S$ collects the source terms (conductivity damping, constraint-damping
relaxation, and material-gradient couplings). The three flux matrices
$A_x, A_y, A_z$ encode the spatial propagation. Their
characteristic decomposition is what the penalty method needs; we return
to this in Sec.~\ref{sec:sat}.

\section{Spatial discretization}
\label{sec:spatial}

\subsection{Grid, decomposition, and indexing}

The domain is a rectangular parallelepiped
$[x_0, x_n] \times [y_0, y_n] \times [z_0, z_n]$, discretized on a
uniform Cartesian grid of global size $(N_x, N_y, N_z)$ with spacings
$dx, dy, dz$. The grid is distributed across MPI ranks via a Cartesian
topology chosen automatically by greedy prime factorization. Each rank
owns a contiguous block and is surrounded by ghost zones of width
$g_s$ (default $g_s = 3$). For physical boundaries the ghost zones are
absent and the rank stores real data at every index; for MPI-internal
and periodic boundaries, the ghost zones are synchronized via
non-blocking point-to-point communication before each RK stage.

The local index into a 3D field $f$ at grid point $(i, j, k)$ is
$\mathtt{ijk} = i + j N_x + k N_x N_y$, with $i$ the fastest-varying
index. Strides are $d_i = 1$, $d_j = N_x$, $d_k = N_x N_y$, and every
finite-difference stencil takes a stride argument so it can operate
along any axis.

\subsection{Summation-by-parts (SBP) operators}
\label{sec:sbp-operators}

The spatial derivatives are computed by summation-by-parts finite
differences of order $(4,2)$ — fourth-order in the interior, second-order
one-sided closures at the first four rows on each boundary. A $(4,2)$
SBP operator $D$ approximates $d/dx$ on a grid of $N$ points and
factorizes as
\begin{equation}
  D = H^{-1} Q, \qquad Q + Q^\top = B = \mathrm{diag}(-1, 0, \ldots, 0, +1),
\end{equation}
where $H$ is a positive-definite diagonal weight matrix (the ``norm''),
and $Q$ is almost antisymmetric — the asymmetry captured entirely by the
boundary term $B$. In the $(4,2)$ case
\begin{equation}
  H = h \cdot \mathrm{diag}(17/48,\, 59/48,\, 43/48,\, 49/48,\, 1, 1, \ldots),
\end{equation}
with $h$ the grid spacing. The key stability-enabling identity is
\begin{equation}
  \langle u, D v \rangle_H + \langle Du, v \rangle_H = u_N v_N - u_0 v_0,
\end{equation}
where $\langle \cdot, \cdot \rangle_H = \sum_i H_{ii} u_i v_i$ is the
discrete $H$-inner product. This is the discrete analogue of integration
by parts on a finite interval — the boundary terms appear exactly,
without truncation error, and it is this identity that lets us prove
a discrete energy estimate for the scheme.

The interior stencil (rows $i = 4 \ldots N-5$) is the standard
fourth-order centered difference
\begin{equation}
  (D u)_i = \frac{u_{i-2} - 8 u_{i-1} + 8 u_{i+1} - u_{i+2}}{12 h}.
\end{equation}
The four near-boundary rows use one-sided Strand (1994) closures;
at the lower boundary:
\begin{align*}
  (D u)_0 &= \frac{-24 u_0 + 59 u_1 - 8 u_2 - 3 u_3}{17 h}, \\
  (D u)_1 &= \frac{-u_0 + u_2}{2 h}, \\
  (D u)_2 &= \frac{ 4 u_0 - 59 u_1 + 59 u_3 - 4 u_4}{86 h}, \\
  (D u)_3 &= \frac{ 3 u_0 - 59 u_2 + 64 u_4 - 8 u_5}{98 h},
\end{align*}
and symmetric closures at the upper boundary. These are implemented as
\texttt{SBP42\_L0}, \texttt{SBP42\_L1}, \texttt{SBP42\_L2},
\texttt{SBP42\_L3} and their upper-boundary counterparts
(\texttt{SBP42\_RN}, etc.) in \texttt{derivatives.h}.

\subsection{Runge--Kutta time integration}

Time integration is classical explicit fourth-order Runge--Kutta. The
timestep is chosen by the Courant--Friedrichs--Lewy condition
\begin{equation}
  dt = \mathrm{cfl} \cdot \min(dx, dy, dz).
\end{equation}
For the pure SBP interior (no SAT), stability permits $\mathrm{cfl}
\lesssim 1$. With SAT penalties on all six faces the effective
spectral radius increases near corners where multiple penalties
stack, and the usable CFL drops to $\approx 0.2$ in practice
(Sec.~\ref{sec:limitations}).

\section{Boundary conditions: the SBP-SAT penalty method}
\label{sec:sat}

\subsection{Why weak imposition}

The na\"ive approach to boundary conditions — ``overwrite the evolution
at boundary nodes with the prescribed boundary value'' — is called
\emph{strong imposition}. It is simple and often works, but it breaks
the SBP energy estimate: a boundary row whose equation is replaced
loses the clean integration-by-parts structure, and stability must be
argued separately, usually by numerical experiment.

The \emph{simultaneous approximation term} (SAT) method imposes boundary
conditions \emph{weakly}, as an additive penalty. The PDE is retained
at every row, including the boundary row, and a correction term is
\emph{added} that pulls the solution toward the prescribed boundary
value. The SAT term is tuned so that the full semi-discrete system
admits a discrete energy estimate: summing the equations against
$U^\top H$ and using the SBP identity produces boundary terms that
are exactly cancelled by the SAT penalty. This gives provable stability
for any $\tau \geq 1/2$ (see below), independent of the particular
PDE coefficients or grid spacing.

The tradeoff is summarised below.

\begin{center}
\begin{tabular}{lll}
\toprule
Property                       & Strong imposition            & SBP-SAT (weak)      \\
\midrule
Boundary equation              & Replaced entirely            & Interior RHS + penalty \\
Energy estimate                & No formal proof              & Yes, provable for $\tau \geq 1/2$ \\
Stability w/ non-uniform $\eps, \mu$ & No guarantee                 & Provable \\
Convergence order              & Same as interior closures    & Same as interior closures \\
Penalty scaling                & N/A                          & $\tau \cdot H^{-1}_{00}/h$ \\
\bottomrule
\end{tabular}
\end{center}

Both methods give the same formal convergence order — the boundary
truncation error is set by the SBP one-sided closure, not by how the
boundary condition is imposed. The pedagogical reason to prefer
SBP-SAT is the discrete energy estimate, which extends to variable
coefficients, multi-dimensional grids, and characteristic couplings
(see Sec.~\ref{sec:stability-tau}) — all places where a strong
imposition would need a separate stability argument.

\subsection{Characteristic decomposition}
\label{sec:characteristics}

The penalty acts on the \emph{incoming} characteristic at each face.
For a hyperbolic system $\partial_t U + A \partial_x U = 0$ with
flux matrix $A$, diagonalize
\begin{equation}
  A = V \Lambda V^{-1}, \qquad \Lambda = \mathrm{diag}(\lambda_1, \ldots, \lambda_n),
\end{equation}
and split into positive and negative eigenvalue parts
\begin{equation}
  A^+ = V \Lambda^+ V^{-1}, \qquad A^- = V \Lambda^- V^{-1},
  \qquad \Lambda^\pm = \mathrm{diag}(\max(\pm \lambda_i, 0)).
\end{equation}
The characteristic variables $W = V^{-1} U$ satisfy $(\partial_t + \lambda_i \partial_x) W_i = 0$;
modes with $\lambda_i > 0$ propagate in $+x$ (incoming at the lower face $x = 0$), modes with
$\lambda_i < 0$ propagate in $-x$ (incoming at the upper face $x = L$).

At the lower-$x$ face, we prescribe the \emph{incoming} characteristics
(the $+\lambda$ modes); the $-\lambda$ modes are outflow and are
determined by the interior dynamics. At the upper-$x$ face, the roles
are reversed.

\subsection{The penalty form}

At the lower-$x$ boundary (grid index $i = 0$), the SAT correction
applied to each variable is
\begin{equation}
  \dot U_0 \;\mathrel{-}=\; \tau \cdot \frac{(H^{-1})_{00}}{h} \cdot A^+ \cdot (U_0 - g),
  \label{eq:sat-lower}
\end{equation}
where $g$ is the prescribed boundary state. The scalar $\tau \geq 1/2$
is the penalty strength; the code uses $\tau = 1$ by default.
$(H^{-1})_{00} = 48/17$ for the $(4,2)$ SBP norm. The stencil operates
only on the single boundary row — rows $i = 1, 2, 3$ receive the SBP
one-sided derivative (Sec.~\ref{sec:sbp-operators}) but \emph{no}
additional SAT term; they contribute to the energy argument through
the norm integral.

The analogous correction at the upper boundary uses $A^-$ (or,
equivalently, the positive-definite projection $|A^-|$ onto modes
with negative eigenvalue) and the incoming data for that face:
\begin{equation}
  \dot U_{N-1} \;\mathrel{-}=\; \tau \cdot \frac{(H^{-1})_{N-1,N-1}}{h}
  \cdot |A^-| \cdot (U_{N-1} - g).
  \label{eq:sat-upper}
\end{equation}
In the $(4,2)$ SBP norm, $(H^{-1})_{N-1,N-1} = (H^{-1})_{00} = 48/17$
by symmetry.

\subsection{Application to Maxwell's equations}

For Maxwell's equations extended with constraint damping, the $z$-axis
flux matrix decomposes into four independent $2 \times 2$ blocks plus
one trivial $1 \times 1$ block:

\begin{center}
\begin{tabular}{llll}
\toprule
Block & Variables & Eigenvalues & Incoming char.\ at lower $z$ ($w^+$) \\
\midrule
EM-1         & $(B_x, D_y)$     & $\pm c$ & $B_x - PP \cdot D_y$ \\
EM-2         & $(B_y, D_x)$     & $\pm c$ & $B_y + PP \cdot D_x$ \\
Div-$B$      & $(B_z, \Psi_B)$  & $\pm 1$ & $B_z - \Psi_B$ \\
Div-$D$      & $(D_z, \Psi_D)$  & $\pm 1$ & $D_z - \Psi_D$ \\
Charge       & $\rho$           & $0$     & (none) \\
\bottomrule
\end{tabular}
\end{center}

Here $c = \sqrt{\eps \mu}^{-1}$ is the local light speed (in our
normalization $c = \sqrt{\mathtt{ieps} \cdot \mathtt{imu}}$) and
$PP = \sqrt{\eps/\mu}^{-1} = \sqrt{\mathtt{ieps}/\mathtt{imu}}$ is the
impedance ratio. The divergence blocks have eigenvalues $\pm 1$ because
constraint violations propagate at the normalized speed set by the
extended system (independent of material). The $w^+$ column lists, for
each block, the left-eigenvector combination that is the incoming
characteristic at the lower-$z$ face; the upper-$z$ face injects the
opposite combination ($w^-$) with the signs of the cross-coupling
terms reversed.

For the EM-1 block $(B_x, D_y)$, the flux matrix in the convention used
by the code is
\begin{equation}
  A_z^{(\mathrm{EM1})} = \begin{pmatrix} 0 & -\mathtt{ieps} \\ -\mathtt{imu} & 0 \end{pmatrix}.
\end{equation}
The right eigenvector for $\lambda = +c$ is $v^+ = (1, -1/PP)^\top$;
the left eigenvector (satisfying $L^+ A = c L^+$) is $L^+ = (1, -PP)^\top$;
and the positive projection is
\begin{equation}
  A^+ = \frac{c}{L^+ \cdot v^+} v^+ (L^+)^\top = \frac{c}{2}
    \begin{pmatrix} 1 & -PP \\ -1/PP & 1 \end{pmatrix}.
\end{equation}
Applying to $(\delta B_x, \delta D_y) = (B_x - g_{B_x}, D_y - g_{D_y})$:
\begin{align}
  (A^+ \delta U)_{B_x} &= \tfrac{c}{2} (\delta B_x - PP \cdot \delta D_y), \\
  (A^+ \delta U)_{D_y} &= \tfrac{c}{2} (-\delta B_x / PP + \delta D_y).
\end{align}
The incoming characteristic variable for this block is
$w^+ = L^+ \cdot U = B_x - PP \cdot D_y$.

\paragraph{Verification on a known solution.} For a $y$-polarised
plane wave travelling in $+z$ (incoming at the lower-$z$ boundary), the
fields satisfy $D_y = \psi$, $B_x = -\psi$ (the sign makes $\vE \times \vB$
point in $+\hz$), so
$w^+ = -\psi - PP \cdot \psi = -(1+PP)\psi \ne 0$ as expected for an
incoming mode — nonzero, so the SAT penalty actually injects. A sign
error that flipped any coupling term (e.g.\ $+PP$ instead of $-PP$)
would give $w^+ = 0$ for this same solution, and the SAT would
silently fail to inject the wave. Analogous derivations for the other
three blocks give the full SAT formula for the lower-$z$ face; the
formulas for the remaining five faces follow by cyclic permutation of
axes and sign reversal of the coupling terms for the upper faces
(Appendix~\ref{app:sat-formulas}).

The implementation lives in \texttt{maxwell\_eqs.c} as six preprocessor
macros: \texttt{APPLY\_SAT\_LOWER\_\{X,Y,Z\}} and
\texttt{APPLY\_SAT\_UPPER\_\{X,Y,Z\}}. Only the lower-$z$ face has
nonzero $g$ (the incoming source field, Sec.~\ref{sec:sources}); the
other five faces use $g = 0$, which corresponds to a perfectly
absorbing boundary for waves leaving the domain.

\subsection{Edges and corners}

At an edge (where two physical faces meet) and at a corner (three
faces), the SAT penalties are \emph{additive}: each face contributes
its own penalty term independently. No cross-term or corner-specific
correction is needed. This additivity is a direct consequence of the
tensor-product structure: the 3D discrete energy norm
$H = H_x \otimes H_y \otimes H_z$ decomposes the boundary contributions
into independent surface integrals over each face, each cancelled by
its own SAT term.

Concretely, at the corner $(i, j, k) = (0, 0, 0)$ with all three axes
physical, the RHS receives
\begin{equation}
  \dot U_{0,0,0} \;\mathrel{-}=\; \tau \left[
    \frac{H^{-1}_{0}}{dx} A_x^+ \delta U +
    \frac{H^{-1}_{0}}{dy} A_y^+ \delta U +
    \frac{H^{-1}_{0}}{dz} A_z^+ \delta U
  \right],
  \label{eq:corner-sat}
\end{equation}
where the three projections act on the \emph{same} state $\delta U$ but
with axis-specific flux matrices. This is the \emph{only} corner-aware
term in the code — the SBP derivatives at the corner use one-sided
closures in each direction independently, and no other correction is
applied.

\subsection{Stability, $\tau$, and CFL at corners}
\label{sec:stability-tau}

The standard SBP-SAT stability theorem (Svärd \& Nordström) requires
$\tau \geq 1/2$ for a \emph{single-face} penalty on a scalar hyperbolic
equation. For a system with $n$ characteristic speeds and a tensor-product
multi-dimensional grid, the same bound applies per face. In the code
we use $\tau = 1$ as a comfortable margin.

At corners, three SAT terms with per-face scaling $\tau (H^{-1}_0)/h$
stack additively on the same boundary row. The effective spectral
radius of the discrete evolution operator grows at the corners, and
the usable CFL number drops. Empirically, the scheme with all six
faces physical is stable at $\mathrm{cfl} \approx 0.2$ with classical
RK4 — tighter than the $\mathrm{cfl} \approx 0.5$ safe for an
interior-only scheme. Raising $\tau$ beyond $1$ tightens the CFL
further; reducing $\tau$ toward $1/2$ relaxes it at the cost of
stability margin.

\section{Source models}
\label{sec:sources}

The lower-$z$ face (or, for the convergence-test source, every
physical face) has nonzero boundary data $g$ corresponding to a
specified incoming electromagnetic field. Three sources are
implemented:
\begin{itemize}
  \item \texttt{plane\_wave}: a bump-enveloped plane wave travelling
    in $+\hz$ (Sec.~\ref{sec:plane-wave}). Pair with periodic $x, y$
    for the formal convergence reference (Option B).
  \item \texttt{gaussian\_beam}: a paraxial collimated beam injected
    at the lower-$z$ face (Sec.~\ref{sec:gaussian-beam}). Used with
    all faces non-periodic for the headline beam-through-lens
    physics demonstration (Option A).
  \item \texttt{te\_waveguide\_mode}: an exact vacuum Maxwell
    TE mode of a rectangular unit-cube waveguide, used as SBP-SAT
    boundary data on every physical face for the four-configuration
    convergence study (Sec.~\ref{sec:convergence}).
\end{itemize}
Selection is via the TOML key \texttt{type} under the \texttt{[source]}
table. Per-source parameters live in dedicated sub-tables
\texttt{[source.plane\_wave]}, \texttt{[source.gaussian\_beam]},
\texttt{[source.te\_waveguide\_mode]}: the driver parses all three
(so switching sources is a one-line \texttt{type} edit), but only the
sub-table matching the active \texttt{type} is read by the analytic
source functions.

\subsection{Plane wave source}
\label{sec:plane-wave}

The plane-wave source is a transverse electromagnetic wave travelling
in $+\hz$, linearly polarised as $\vE = (E_0^x \hx + E_0^y \hy)
f(t - z) \cos(k(z - t))$, modulated by a compactly supported smooth
envelope $f$ that turns on over the interval $[a, b]$. In the code
this is \texttt{incoming\_plane\_wave}; the envelope is a $C^2$
polynomial bump. The wave satisfies the vacuum Maxwell equations
exactly, independent of wavenumber $k$, so in a configuration with
$x, y$ periodic (and $z$ non-periodic) this source provides an
\emph{exact} reference solution for formal convergence testing of
the SBP-SAT scheme (Sec.~\ref{sec:using}).

The plane-wave source has no transverse ($x, y$) structure —
the field at $x = 0$ and $x = L_x$ is identical. Combining this source
with non-periodic $x, y$ boundaries therefore produces spurious
reflections: the full-amplitude wavefield at the transverse walls is
pulled toward zero by the $g = 0$ SAT there, exciting high-frequency
numerical modes and ultimately destabilizing the simulation. The plane
wave is not usable with six-face SAT.

\subsection{Paraxial Gaussian beam}
\label{sec:gaussian-beam}

The Gaussian-beam source models a collimated laser beam entering at the
lower-$z$ face, optionally converging toward a waist at $z = z_w$. In
the paraxial approximation, the scalar amplitude is
\begin{equation}
  \psi(\vx, t) = E_0 \cdot \frac{w_0}{w(z)} \exp\!\left( -\frac{r^2}{w(z)^2} \right)
    \cos\!\left( k z - \omega t + \frac{k r^2}{2 R(z)} - \varphi_G(z) \right),
  \label{eq:gauss-amp}
\end{equation}
where $r^2 = x^2 + y^2$, $\omega = k c$ (with $c = 1$ in vacuum units),
and
\begin{align}
  z_R &= \tfrac{1}{2} k w_0^2 \qquad \text{(Rayleigh range)} \\
  w(z)       &= w_0 \sqrt{1 + ((z - z_w)/z_R)^2} \\
  \frac{1}{R(z)} &= \frac{z - z_w}{(z - z_w)^2 + z_R^2} \\
  \varphi_G(z) &= \arctan((z - z_w)/z_R) \qquad \text{(Gouy phase)}.
\end{align}
The field is taken linearly polarised along $\hx$ to leading paraxial
order: $D_x = \psi$, $B_y = \psi$, all other components zero.

\paragraph{Convergence / divergence.} The sign of $1/R(z)$ determines
whether the wavefront is converging or diverging. For $z < z_w$,
$R < 0$: the wavefronts are concave toward the beam axis, the beam
\emph{converges} toward the waist. For $z > z_w$, $R > 0$: the beam
diverges. The beam width at the entrance face $z = 0$ is
\begin{equation}
  w(0) = w_0 \sqrt{1 + (z_w/z_R)^2},
\end{equation}
which equals $w_0$ when $z_w = 0$ (beam at its natural focus at the
entrance) and grows with $z_w$. Placing $z_w$ in the middle of the
domain produces a beam that converges through the first half, reaches
minimum spot at the middle, then diverges.

\paragraph{Collimation.} A beam is ``nearly collimated'' (laser-like)
over a propagation distance $L$ when $z_R \gg L$. The spot size varies
by
\begin{equation}
  \frac{w(L)}{w_0} = \sqrt{1 + (L/z_R)^2},
\end{equation}
so $z_R = 10 L$ gives 0.5\% broadening, and $z_R = 3 L$ gives 5\%.
For fixed domain $L$ this requires large $k w_0^2$, which costs
transverse resolution: see Sec.~\ref{sec:limitations}.

\paragraph{Paraxial residual.} Expression \eqref{eq:gauss-amp} satisfies
the full Maxwell equations only to $O(1/(k w_0)^2)$. The residual
manifests as small but nonzero $\nabla \cdot \vE$ and $\nabla \cdot \vB$
when the beam is injected. The constraint-damping fields $\Psi_D$ and
$\Psi_B$ absorb and radiate away this residual on timescale
$1/\kappa_{D,B}$. A sufficiently paraxial beam ($k w_0 \gtrsim 5$)
leaves a residual below one part in $25$, well within the constraint-
damping system's domain of validity.

\subsection{Smooth turn-on in time}

At $t = 0$, the full-domain initial data is filled with the analytic
source evaluated at $t = 0$. For the plane-wave source, the envelope
$f(t - z)$ vanishes at $t = 0$ for all $z \geq 0$, so the initial state
is identically zero and the wave enters gradually through the SAT
boundary.

For the Gaussian beam, the analytic formula \eqref{eq:gauss-amp} does
not vanish at $t = 0$: filling the interior with it would seed the
domain with a state that is only a paraxial approximation to a
discrete solution. Empirically, this excites a numerical instability
that grows exponentially and overwhelms the simulation on a timescale
of a few light-crossing times, regardless of CFL. To avoid this, the
beam source is multiplied by a smooth $C^2$ turn-on envelope
$f(t, t_a, t_b)$ that rises from $0$ at $t \leq t_a$ to $1$ at
$t \geq t_b$. The turn-on interval $[t_a, t_b]$ is exposed as
\texttt{ramp\_a} and \texttt{ramp\_b} under \texttt{[source.gaussian\_beam]};
the default is $[0, 1]$. The initial state is therefore identically
zero, and the beam grows in through the boundary SAT as the envelope
rises. After $t \geq t_b$ the source is at full amplitude and the
simulation runs in its intended regime. Lengthening the ramp (e.g.\
$t_b = 3$) reduces the paraxial-residual excitation during turn-on,
at the cost of delaying the onset of interesting physics; this is
useful when investigating the residual at high $k \cdot w_0$.

The \texttt{te\_waveguide\_mode} source is an exact solution of the
continuum Maxwell PDE and therefore needs no turn-on — it is
non-zero at $t = 0$, and the initial state on the full grid is
consistent with the boundary data at every face. The convergence-
test routine relies on this consistency (Sec.~\ref{sec:convergence}).

\section{Limitations}
\label{sec:limitations}

\subsection{Transverse-face reflections}

The $g = 0$ SAT on the $\pm x$ and $\pm y$ faces is correct only insofar
as the physical solution has negligible amplitude at those walls. For a
Gaussian beam with $w_0 \ll L_x, L_y$ the tail at the wall decays as
$\exp(-(L_{x,y} / 2 w)^2)$ and reflections are below roundoff for
$L_{x,y} \gtrsim 10 w_0$. For a plane wave (which has full amplitude at
the walls) this boundary treatment is not appropriate — the spurious
reflection is large and cascades into an instability.

\subsection{Paraxial approximation}

The Gaussian-beam source is a paraxial approximation to Maxwell's
equations. The residual constraint violation is $O(1/(k w_0)^2)$ and is
absorbed by the constraint-damping fields. This is fine as a physics
source — real laser beams are described by the same paraxial formula —
but means the beam is not an exact solution of the full Maxwell system.
The $L^2$ error of the numerical evolution against the paraxial
reference is therefore not a convergence metric for the scheme; use the
plane-wave source for that.

\subsection{Resolution / containment / paraxial tradeoff}

Three constraints on the beam parameters:
\begin{itemize}
  \item \textbf{Grid resolution:} $k \cdot dx \lesssim 0.6$, i.e.\ roughly
        10 points per wavelength.
  \item \textbf{Paraxial validity:} $k w_0 \gtrsim 5$, giving residual
        $\leq 4\%$.
  \item \textbf{Transverse containment:} $L_x / w_0 \gtrsim 5$
        (stronger: $\geq 10$) so the Gaussian tail is negligible at
        the $x, y$ walls.
\end{itemize}
For a laser-like (nearly-collimated) beam, the Rayleigh range
$z_R = k w_0^2 / 2$ must additionally be large compared with the
propagation length $L_z$. All four constraints together favour large
$k$ and large $w_0$, which raises the required $N_x$ steeply. Practical
configurations balance collimation against computational cost.

\subsection{CFL at corners}

With all six faces treated by SAT, the scheme is stable at
$\mathrm{cfl} \approx 0.2$ with classical RK4 — tighter than the
$\mathrm{cfl} \approx 0.5$ safe for the periodic-xy interior-only
scheme. This is a consequence of the three SAT penalties stacking at
corners (Sec.~\ref{sec:sat}). If you see exponentially-growing $L^2$
error, reduce the CFL before suspecting the physics.

\subsection{Plane wave with six-face SAT}

The plane-wave source combined with all six faces set to SAT produces a
numerical instability (the transverse walls see full-amplitude
tangential field and the $g = 0$ SAT attempts to force it to zero). The
plane-wave source is supported only with $x, y$ periodic and $z$
non-periodic. Attempting the six-face-SAT plane-wave run is \emph{not}
a bug; it is the documented failure mode of the SAT boundary for a
physical configuration whose field is not well-contained in the
transverse plane.

\subsection{Angular (oblique) reflections from every SAT face}
\label{sec:oblique-reflection}

The most important limitation of the current boundary treatment is one
that all characteristic-based absorbing boundaries share: the SAT
penalty is exact for waves moving \emph{normal} to the boundary, but
only partial for waves moving at oblique angles. This affects every
face, not just the transverse ones:

\begin{itemize}
  \item \textbf{Lower-z face}: the SAT's $+c$ projection correctly
        leaves a head-on outgoing wave ($-c$ mode) untouched. A wave
        moving at angle $\theta$ to the $z$-axis has its energy split
        between the $+c$ and $-c$ z-characteristics in proportions that
        depend on $\sin\theta$. The SAT treats the $+c$ part as
        incoming and pulls it toward the prescribed beam data,
        producing a spurious reflection whose amplitude grows with
        $\theta$.
  \item \textbf{Upper-z face}: same mechanism, with $g = 0$. A wave
        refracted by the lens and travelling mostly in $+z$ but at an
        angle is partially reflected back into the domain.
  \item \textbf{$\pm x$, $\pm y$ faces}: same mechanism yet again. The
        beam reaches these walls only as a Gaussian tail and as
        lens-scattered components, so their contribution is small for
        well-contained beams but non-negligible whenever the lens
        refracts energy sideways.
\end{itemize}

\paragraph{Concrete failure mode: beam on a high-$\varepsilon$ lens.}
A dielectric step of permittivity $\varepsilon_{\max}$ reflects a
Fresnel fraction
$R = \left[(\sqrt{\varepsilon_{\max}} - 1)/(\sqrt{\varepsilon_{\max}}
+ 1)\right]^2$ of incident intensity at normal incidence. For
$\varepsilon_{\max} = 4$, $R \approx 0.11$ in intensity, or $1/3$ in
amplitude — a substantial portion of the incident beam bounces off the
lens and propagates back toward $-z$. The pure-$-z$ component exits
cleanly through the lower-z SAT; components scattered at angles by the
elliptical side-lobes of the lens are partially reflected and start a
second pass through the lens. After a few such bounces the domain
fills with low-amplitude oblique modes.

\paragraph{Observable symptoms.} High-frequency interference patterns
(the incident + reflected beam forming standing-wave nodes every
$\lambda/2 = \pi/k$) and visible oblique wavefronts near the lower-z
face at late times. The interference is real physics; the oblique
wavefronts arriving at the lower-z face are the spurious angular
reflections. Both coexist in any long-time run with a moderately
strong lens.

\paragraph{Mitigations within the current code.}
\begin{itemize}
  \item \textbf{Shorten the simulation time} so no wave has time to
        complete a round trip.
  \item \textbf{Weaken the lens} ($\varepsilon_{\max}$ closer to 1) so
        the reflection coefficient is smaller and the artifacts drop
        proportionally.
  \item \textbf{Widen the domain} so the transverse walls are far
        from the beam even after the lens refracts it.
  \item \textbf{Enable Kreiss-Oliger dissipation}
        (\texttt{use\_dissipation = true}, \texttt{diss\_coeff} $\sim$
        $0.02$--$0.1$) to damp grid-scale noise from partial
        reflections. This does not remove the angular reflections
        themselves but suppresses their high-frequency signature.
  \item \textbf{Do not use the current code for
        long-time resonance studies}: the SAT is not designed to model
        a cavity with clean walls.
\end{itemize}

None of these mitigations is a true fix. To absorb oblique waves at all
angles, a different class of boundary treatment is required —
the Perfectly Matched Layer (next section).

\section{Perfectly Matched Layers: a future extension}
\label{sec:pml}

A Perfectly Matched Layer (PML) is a finite-thickness absorbing buffer
that surrounds the physical domain and is designed to absorb outgoing
waves \emph{at any angle and any frequency}, with zero reflection at
the interface with the physical domain at the continuum level. PML
is the standard tool in production electromagnetic and elastodynamic
codes for open-domain simulations. The current \texttt{Maxwell\_Penalty}
code does not include a PML — the SBP-SAT boundary is the only
absorbing mechanism — and this section is a forward-looking note on
what a PML implementation would add and how it would fit.

\subsection{The physical picture}

The idea, due to B\'erenger~(1994), is to surround the computational
domain with a thin layer (typically 8--20 cells) in which the
governing equations are modified so that:

\begin{enumerate}
  \item A wave crossing the physical-domain / PML interface sees no
        discontinuity: its wavenumber, direction, and phase velocity
        are unchanged. No reflection is generated at the interface
        — this is the ``perfectly matched'' property.
  \item Once inside the layer, the wave decays exponentially as it
        propagates away from the physical domain.
  \item Any residual wave that reaches the outer face of the PML
        and reflects (from whatever boundary condition is imposed
        there — a trivial zero Dirichlet suffices) is further
        attenuated on the return trip through the layer.
\end{enumerate}

The round-trip attenuation through the PML controls the effective
reflection the physical domain sees. With a $\sim 10$-cell layer and
a standard $\sigma$ profile, round-trip amplitude ratios of $10^{-4}$
to $10^{-6}$ are routine — orders of magnitude better than the
characteristic-based SAT for oblique incidence.

\subsection{The mathematical mechanism: complex coordinate stretching}

For Maxwell's equations in, say, an $x$-face PML, one formally
replaces the $x$-derivatives by stretched derivatives
\begin{equation}
  \frac{\partial}{\partial x} \;\longrightarrow\;
  \frac{1}{s_x(x)} \frac{\partial}{\partial x},
  \qquad
  s_x(x) \;=\; 1 + \frac{\sigma_x(x)}{i \omega}.
  \label{eq:pml-stretch}
\end{equation}
The stretch factor $s_x$ is complex; $\sigma_x(x) \geq 0$ is the
``PML damping profile'', equal to zero in the physical domain and
ramping up through the PML region.

Substituting a plane wave $e^{i(k x - \omega t)}$ into the stretched
equation gives, in the PML,
\[
  e^{i (k x - \omega t)} \cdot \exp\!\left( -\int_{x_0}^{x}
  \frac{\sigma_x(x')}{c} \, dx' \right),
\]
i.e.\ the same wave with an exponential envelope. Crucially, the
wavenumber $k$ and the phase velocity are unchanged, so the wave does
not refract at the PML interface — which is why the match is perfect
for \emph{all} angles and \emph{all} frequencies.

Similar stretch factors apply at the $y$ and $z$ faces. In a corner
region where two or three PMLs overlap, all the corresponding
derivatives are stretched; the continuum analysis still gives no
corner reflection.

\subsection{From frequency domain to time domain}

The complex factor $1/s_x = 1 / (1 + \sigma_x/(i\omega))$ is trivial
in the frequency domain but requires care in the time domain —
dividing by $(i\omega + \sigma_x)$ is a convolution in time with an
exponential kernel. Two standard implementations:

\begin{itemize}
  \item \textbf{B\'erenger split-field PML (1994).}  Each field
        component is split into pieces labelled by the derivative
        that contributes to it; each piece gets its own damping
        equation. Simple but doubles the number of field variables
        in the PML region.
  \item \textbf{Convolutional PML / CPML (Roden~\&~Gedney 2000).}
        Uses the same complex stretch but implements the convolution
        directly through one auxiliary field per stretched
        derivative, evolved by a first-order ODE. Far fewer extra
        fields, more robust at grazing incidence, and handles
        inhomogeneous media (like our lens) more gracefully. CPML is
        the modern standard and would be the recommended choice
        here.
\end{itemize}

In CPML, each stretched derivative
$(1/s_x) \partial_x E$ in the time-domain PDE is replaced by
$\partial_x E + \Psi^x$, where $\Psi^x$ is an auxiliary field
evolved by
\[
  \frac{d \Psi^x}{dt} \;=\; -\alpha_x \Psi^x \;-\; \sigma_x \,
  \partial_x E,
\]
with $\alpha_x$ an optional additional parameter (often zero) that
improves performance at grazing incidence. Outside the PML,
$\sigma_x = \alpha_x = 0$ and $\Psi^x$ is identically zero; inside,
$\Psi^x$ tracks the history of the spatial derivative damped
exponentially.

\subsection{The damping profile}

$\sigma(x)$ must start at zero at the PML / physical-domain
interface and rise smoothly to a maximum at the outer face:
\begin{equation}
  \sigma(x) \;=\; \sigma_{\max}
  \left( \frac{x - x_{\text{interface}}}{L_{\text{PML}}} \right)^m,
  \qquad m = 2, 3, 4 \text{ typical},
\end{equation}
with $m = 3$ (cubic profile) a standard choice. The maximum
$\sigma_{\max}$ is tuned so the theoretical round-trip reflection
\begin{equation}
  R \;=\; \exp\!\left(
    -\frac{2 \sigma_{\max} L_{\text{PML}}}{(m+1) \, c}
  \right)
\end{equation}
sits at the target level (typically $R \sim 10^{-5}$). Starting
$\sigma$ exactly at zero at the interface is essential: any
discontinuity in $\sigma$ at that boundary creates a small but
visible reflection, and the whole point of a PML is to avoid it.

\subsection{What a PML would add to \texttt{Maxwell\_Penalty}}

Adding CPML to the six faces would require, roughly:

\begin{enumerate}
  \item \textbf{Grid extension.} Physical domain padded by
        $N_{\text{PML}} \approx 10$ cells on each physical face, so
        $N_x \to N_x + 2 N_{\text{PML}}$ etc. A moderate overhead;
        for $N = 64$, a 10-cell PML per face grows the grid to $84$
        per axis, so the total work grows by
        $(84/64)^3 \approx 2.3\times$.
  \item \textbf{Auxiliary fields.} In CPML one needs six
        $\Psi^{\text{axis}}_{\text{component}}$ fields (one per
        stretched derivative per field component), evolved on the
        PML region only. Memory overhead is modest because $\Psi$ is
        zero outside the PML; storing it only in the layer keeps
        the overhead to $\sim 30\%$ of the PML-region cost.
  \item \textbf{RHS extension.} The \texttt{SIMPLE\_MAXWELL\_INTERIOR\_DOT}
        macro would need a PML-aware cousin that includes the $\Psi$
        contributions and evolves them. Because $\sigma$ varies with
        position, the PML body is not easily macro-ised with fixed
        coefficients; a small per-point helper function is more
        natural.
  \item \textbf{Constraint-damping coverage.} The extended Maxwell
        system here carries $\Psi_D$ and $\Psi_B$ with their own
        constraint-propagation speeds (normalized to $1$, not $c$).
        These blocks need their own PML damping profile, usually
        tuned to match the same round-trip attenuation target. The
        constraint-damping and EM PMLs have different physical
        wavelengths so their $\sigma_{\max}$ values differ.
  \item \textbf{Outer boundary.} The outermost cell of the PML can
        use a trivial zero-Dirichlet boundary. A zero SAT would also
        work. The round-trip attenuation makes the choice
        effectively irrelevant.
  \item \textbf{MPI topology.} The PML adds layer cells on ranks
        that own physical faces, but the MPI decomposition is
        otherwise unchanged. Ghost-zone synchronisation on the
        interior MPI faces works as before.
\end{enumerate}

The complete implementation effort is on the order of a few hundred
lines of code and a week of careful convergence testing. The payoff
is clean simulations of open-domain scattering problems — a lens
demo, a beam refraction study, a waveguide mode leakage analysis —
with artifacts driven below the single-precision noise floor.

\subsection{Summary: SAT versus PML}

\begin{center}
\begin{tabular}{lll}
\toprule
 & SBP-SAT (current) & PML (future extension) \\
\midrule
Mechanism          & Characteristic projection & Complex coordinate stretching \\
Normal incidence   & Exact (for $\tau \geq 1/2$) & Near-exact, $R \sim 10^{-5}$ \\
Grazing incidence  & Strongly reflective        & Near-exact \\
Extra grid cells   & None                       & $\sim 10$ per physical face \\
Extra fields       & None                       & One $\Psi$ per stretched derivative \\
Energy estimate    & Yes (provable)             & Not directly; effective via round-trip bound \\
Handles lens refl. & Head-on only               & All angles \\
Code complexity    & Low                        & Medium--high \\
\bottomrule
\end{tabular}
\end{center}

The choice between SAT and PML is a design decision, not a
correctness one: both can be implemented correctly; they just
handle different regimes well. For a teaching code demonstrating
SBP-SAT, SAT alone is appropriate. For production scattering
simulations, PML is the right tool.

\section{Using the code}
\label{sec:using}

\subsection{Build}

Dependencies: MPI (tested with OpenMPI), HDF5 (with the high-level
interface), a C11 compiler, and a C++17 compiler (for the TOML
parameter parser).

From \texttt{src/}:
\begin{lstlisting}[language=bash]
cmake -S . -B build
cmake --build build -j
\end{lstlisting}
or equivalently with the hand-written Makefile:
\begin{lstlisting}[language=bash]
make maxwell_system
\end{lstlisting}
The resulting executable is \texttt{maxwell\_system}. A test suite
(domain decomposition, ghost-zone synchronization, RK4 convergence) is
run with
\begin{lstlisting}[language=bash]
ctest --test-dir build
\end{lstlisting}

\subsection{Running}

\begin{lstlisting}[language=bash]
mpirun -np <N> ./maxwell_system maxwell.toml
\end{lstlisting}
Output is written to \texttt{3D\_rank\_<R>.h5} for each rank
(one file per MPI rank, one group per output step). See
Sec.~\ref{sec:visualization} for how to visualize these files.

\subsection{Parameter file structure}

The TOML parameter file has the following sections:

\paragraph{[grid]} Domain extent and resolution.
\begin{lstlisting}[language={}]
nx = 64          # global grid points in x (total, not per rank)
ny = 64
nz = 128
x0 = -1.0        # lower corner
y0 = -1.0
z0 =  0.0
xn =  1.0        # upper corner
yn =  1.0
zn =  4.0
periodic_x = false  # if true, no SAT on this axis, ghost zones wrap
periodic_y = false
periodic_z = false
\end{lstlisting}

\paragraph{[solver]} Numerical parameters.
\begin{lstlisting}[language={}]
ghost_size       = 3
cfl_factor       = 0.2    # use <= 0.2 for all-physical; <= 0.5 for periodic x,y
max_iterations   = 2049
output_every     = 32
checkpoint_every = 512
max_checkpoints  = 2
recover          = false
use_dissipation  = false  # Kreiss-Oliger, off by default
diss_coeff       = 0.1
tau              = 1.0    # SAT penalty strength, >= 1/2
\end{lstlisting}

\paragraph{[physics]} Constraint-damping rates.
\begin{lstlisting}[language={}]
kappa_D = 0.1
kappa_B = 0.1
\end{lstlisting}

\paragraph{[source]} Selects the incoming source type. Per-source
parameters live in dedicated sub-tables, of which exactly one is
consumed per run (matching the \texttt{type} value); the others are
parsed but ignored, so switching sources is a one-line \texttt{type}
edit with no other TOML shuffling.
\begin{lstlisting}[language={}]
[source]
type = "gaussian_beam"   # or "plane_wave" or "te_waveguide_mode"
\end{lstlisting}

\paragraph{[source.plane\_wave]} Plane-wave parameters. The envelope
turn-on interval $[a, b]$ is in $t - z$ (characteristic) coordinate.
\begin{lstlisting}[language={}]
[source.plane_wave]
ax     = 1.0       # x-polarisation amplitude
ay     = 1.0       # y-polarisation amplitude
k      = 4.0       # wavenumber
bump_a = 0.0       # envelope turn-on start
bump_b = 2.0       # envelope turn-on end
\end{lstlisting}

\paragraph{[source.gaussian\_beam]} Paraxial Gaussian-beam parameters.
The temporal ramp \texttt{[ramp\_a, ramp\_b]} is the $[t_a, t_b]$
interval of the smooth $C^2$ turn-on envelope
(Sec.~\ref{sec:gaussian-beam}); lengthen it to reduce the paraxial-
residual excitation during turn-on.
\begin{lstlisting}[language={}]
[source.gaussian_beam]
w0        = 0.20   # waist radius
z_waist   = 2.0    # z location of waist (inside domain for lens demo)
k         = 15.0   # wavenumber
amplitude = 1.0    # peak field amplitude
ramp_a    = 0.0    # turn-on start
ramp_b    = 1.0    # turn-on end
\end{lstlisting}

\paragraph{[source.te\_waveguide\_mode]} Rectangular waveguide TE
mode numbers. \texttt{l, m} are the transverse standing-wave orders
in $x, y$; \texttt{n} is the longitudinal (propagation-direction)
order in $z$. All three must be $\geq 1$. Used for the convergence
study (Sec.~\ref{sec:convergence}).
\begin{lstlisting}[language={}]
[source.te_waveguide_mode]
l = 2
m = 2
n = 2
\end{lstlisting}

\paragraph{[material]} The selector \texttt{epsilon\_type} chooses how
permittivity varies in space. \texttt{[material.background]} carries
the always-used uniform scalars (permittivity, permeability,
conductivity); when \texttt{epsilon\_type = "elliptical"},
\texttt{[material.elliptical]} additionally describes a lens-shaped
bump layered on top of the background via
\[
  \varepsilon(x, z) =
    (\varepsilon_\text{max} - \varepsilon_\text{bg})\,
      \exp\!\left[-s^2 \left(\frac{(x - x_0)^2}{a^2} + \frac{(z - z_0)^2}{b^2}\right)^2\right]
    + \varepsilon_\text{bg},
\]
where $\varepsilon_\text{bg}$ is the \texttt{epsilon} value from
\texttt{[material.background]}. With the default
$\varepsilon_\text{bg} = 1$ this reduces to the vacuum-background
lens used for the headline physics demo.
\begin{lstlisting}[language={}]
[material]
epsilon_type = "elliptical"   # or "uniform"

[material.background]
epsilon = 1.0                 # baseline permittivity (asymptote of the elliptical profile)
mu      = 1.0
sigma   = 0.0

[material.elliptical]
max = 4.0                     # peak permittivity at the lens centre
x0  = 0.0                     # centre x
z0  = 2.0                     # centre z
s   = 2.0                     # steepness
a   = 0.5                     # semi-axis x
b   = 0.25                    # semi-axis z
\end{lstlisting}

\subsection{Two canonical configurations}

\paragraph{Option B: plane-wave convergence reference.}
Periodic $x, y$, non-periodic $z$, plane-wave source. The only
configuration in which an exact, non-dispersive travelling-wave
solution exists for the plane-wave source (non-periodic transverse
faces would force a $g = 0$ SAT at full-amplitude wavefield points
and destabilise the run; see Sec.~\ref{sec:limitations}):
\begin{lstlisting}[language={}]
[grid]
periodic_x = true
periodic_y = true
periodic_z = false

[source]
type = "plane_wave"
\end{lstlisting}
On a sequence of refined grids, the $L^2$ error relative to the
analytic solution reduces at the order of the SBP-$(4,2)$ boundary
closure (formally second-order locally at the boundary, integrating
to higher order globally; empirically around third-order).

\paragraph{Option A: beam through lens.}
All six faces non-periodic with SAT, Gaussian-beam source,
elliptical permittivity. This is the ``headline'' physics
demonstration: a collimated (or mildly convergent) laser beam
enters at $z = 0$, propagates through a lens-shaped high-
permittivity region at $z \approx z_w$, and exits at $z = L_z$.
\begin{lstlisting}[language={}]
[grid]
periodic_x = false
periodic_y = false
periodic_z = false

[source]
type = "gaussian_beam"

[material]
epsilon_type = "elliptical"
\end{lstlisting}

\subsection{Output layout}

Each rank writes a self-describing HDF5 file
\texttt{3D\_rank\_<R>.h5} containing the full field state at every
output interval, grouped by output index. Groups $/0, /1, /2, \ldots$
correspond to the initial state, the first output step, and so on
(step $N$ is at simulation time
$N \cdot \mathtt{output\_every} \cdot dt$). Each group has datasets
for the nine evolved variables and the five auxiliary variables
(including the diagnostic divergences $c_D$, $c_B$). The \texttt{/metadata}
group carries per-rank attributes describing the grid (origin, spacing,
local and global sizes, ghost width, variable names) — everything
needed to reassemble the global domain without any external bookkeeping.

Note that with the beam source's smooth turn-on, group $/0$ contains
identically zero fields. The first non-trivial output is at
$t = \mathtt{output\_every} \cdot dt$; by $t \geq \mathtt{ramp\_b}$
the turn-on envelope has saturated and the beam is at full amplitude.

An $L^2$ diagnostic \texttt{l2\_norm.dat} is written by rank 0 every
output step, containing $(t, \Vert U_\text{numerical} - U_\text{analytic} \Vert_{L^2})$.
For the plane-wave source with periodic $x, y$ this is a convergence
metric; for the Gaussian-beam source the analytic reference is only the
paraxial approximation and the $L^2$ error reaches a paraxial-scale
plateau, not a convergence-to-zero.

\subsection{Visualization}
\label{sec:visualization}

The raw HDF5 layout is not directly readable by ParaView, VisIt, or
PyVista — those tools expect either a VTK-native format or a sidecar
manifest that describes the grid. The preferred workflow is to emit an
XDMF sidecar after the run:

\begin{lstlisting}[language=bash]
python scripts/make_xdmf.py --dir <run_dir> --dt-per-output <T>
\end{lstlisting}

where \texttt{<T>} is the physical time between successive HDF5 output
groups, i.e.\
$T = \mathtt{output\_every} \cdot \mathtt{cfl\_factor} \cdot \min(dx, dy, dz)$.
The script scans the per-rank HDF5 files, reads their \texttt{/metadata},
and emits a single \texttt{maxwell.xmf} manifest alongside the HDF5
files. The manifest is pure XML metadata — the payload bytes stay in
HDF5 and are never copied.

The resulting \texttt{.xmf} file is a direct input for all three
target tools:
\begin{itemize}
  \item \textbf{ParaView}: \texttt{File $\to$ Open $\to$ maxwell.xmf}.
    The Xdmf3 reader is preferred; the legacy Xdmf2 reader also
    works.
  \item \textbf{VisIt}: \texttt{File $\to$ Open $\to$ maxwell.xmf}.
    VisIt's Xdmf plugin reads it directly.
  \item \textbf{PyVista}: \texttt{pyvista.XdmfReader("maxwell.xmf")}.
\end{itemize}

A time slider stepping through the $N$ output steps appears
automatically; all 14 fields are selectable as point-data attributes.

\paragraph{Grid-layout details.} The XDMF uses a 3DRectMesh topology
with an explicit VXVYVZ geometry (one per-axis coordinate list).
This avoids the long-standing XDMF ORIGIN\_DXDYDZ ambiguity whereby
VisIt and ParaView disagree on whether the origin tuple is listed in
spatial $(x, y, z)$ order or in topology-dimensions $(z, y, x)$ order:
VXVYVZ tags each coordinate list explicitly with its axis role, so
both readers agree.

Each MPI rank appears as a spatial sub-grid inside each time step.
Adjacent ranks share one grid plane at each internal MPI boundary
(one extra ghost row is included on the upper side of every MPI
boundary), so the rendered cells tile the global grid without gaps.
The shared points are filled from the same ghost-exchanged data
before output, so the two neighbours render identical values at the
shared plane and the overlap is invisible.

\paragraph{Legacy path.} The older \texttt{scripts/hdf5\_to\_vtk.py}
reassembles the global grid on rank 0 and writes one VTI per time
step. It is still supplied as a fallback for environments where XDMF
is not convenient, but it is serial, duplicates the data on disk, and
has historically surfaced platform-specific portability bugs (e.g., a
NumPy 2.0-era issue in which \texttt{str(np.float64(x))} became
\texttt{"np.float64(x)"} and leaked into the VTI XML header; if you
have pre-existing VTI files with that issue, a tiny in-place patcher
\texttt{scripts/fix\_vti\_header.py} corrects the header without
rewriting the data payload).

\subsection{Diagnostics and common problems}

\paragraph{$L^2$ error blows up exponentially.} Usually CFL is too
large. Reduce \texttt{cfl\_factor} to $0.2$ (six-face SAT) or verify
the configuration is in fact Option B if you expected $\mathrm{cfl} = 0.4$
to work.

\paragraph{$D$ field looks ``empty'' in visualization.} Check the time
slice. Group $/0$ is the initial state, which is zero for both the
plane-wave and Gaussian-beam sources (the turn-on envelope). For
the default ramp $[0, 1]$ and the default \texttt{output\_every} on
the shipped TOML, the beam has turned on to $\sim 3\%$ by group
$/1$ ($t = 0.2$) and reaches full amplitude by $t \geq \mathtt{ramp\_b}$;
visualize group $/5$ or later for the default parameters. (The
\texttt{te\_waveguide\_mode} source has no ramp and its group $/0$
is already the full analytic field.)

\paragraph{Paraxial residual looks larger than expected.} Verify
$k \cdot w_0 \gtrsim 5$. Residual scales as $1/(k w_0)^2$.

\paragraph{Beam appears to ``reflect'' off the transverse walls.}
The $x, y$ faces use $g = 0$ SAT; the reflection amplitude is
proportional to the field amplitude at the wall. Increase $L_x$,
$L_y$, or reduce $w_0$ so the tail at the wall is negligible. Rule
of thumb: $L_x \geq 10 w_0$ where $w_0$ here should be the local spot
size along the propagation, not just the waist.

\paragraph{At late times, high-frequency structure and oblique
wavefronts appear near the lower-z face.} This is the combined
signature of a real standing wave (incident + lens-reflected
beam interfering every $\lambda/2 = \pi/k$) and spurious
angular reflections of lens-scattered energy off every SAT face.
The first is physics; the second is a known limitation of the SAT
boundary treatment — see Sec.~\ref{sec:oblique-reflection}. Enabling
Kreiss-Oliger dissipation (\texttt{use\_dissipation = true},
\texttt{diss\_coeff} $\sim 0.05$) suppresses the high-frequency
component. For truly angle-independent absorption the code would
need a PML — see Sec.~\ref{sec:pml}.

\section{Verification: convergence testing}
\label{sec:convergence}

A direct test for the correctness of the interior stencil, each
face's SBP-SAT closure, and the additive penalty at edges and corners
is obtained by forcing the numerical solution to converge to a known
analytic solution of the continuum PDE and measuring the convergence
rate. This section describes the test setup implemented in the code
and how to run it.

\subsection{Test problem}

The test uses the \texttt{te\_waveguide\_mode} function from
\texttt{analytic\_solutions.c}, which evaluates an exact vacuum
Maxwell solution on $[0, 1]^3$ with uniform permittivity
$\varepsilon = \mu = 1$, zero conductivity, and no sources. The
solution is the TE rectangular-waveguide mode with transverse
standing-wave numbers $(l, m)$ in $(x, y)$ and longitudinal
propagation mode $n$ in $z$:
\begin{align*}
  D_x(x, y, z, t) &= \frac{\omega}{l \pi}\, \sin(m \pi y)\, \cos(l \pi x)\,
                     \cos(n \pi z - \omega t), \\
  D_y(x, y, z, t) &= -\frac{\omega}{m \pi}\, \sin(l \pi x)\, \cos(m \pi y)\,
                     \cos(n \pi z - \omega t), \\
  D_z(x, y, z, t) &= 0, \\
  B_x(x, y, z, t) &= \frac{n}{m}\, \sin(l \pi x)\, \cos(m \pi y)\,
                     \cos(n \pi z - \omega t), \\
  B_y(x, y, z, t) &= \frac{n}{l}\, \sin(m \pi y)\, \cos(l \pi x)\,
                     \cos(n \pi z - \omega t), \\
  B_z(x, y, z, t) &= -\frac{l^2 + m^2}{l m}\, \sin(n \pi z - \omega t)\,
                     \cos(l \pi x)\, \cos(m \pi y),
\end{align*}
with constraint fields $\Psi_D = \Psi_B = \rho = 0$ and the
dispersion relation $\omega = \pi \sqrt{l^2 + m^2 + n^2}$. The
$l, m, n$ mode numbers are set under \texttt{[source.te\_waveguide\_mode]}
in the TOML (all three integer, $\geq 1$; the parameter parser
enforces the positivity bound, which also rules out the $1/(lm)$
division below). The shipped default is $(l, m, n) = (2, 2, 2)$.
The solution satisfies $\nabla \cdot \vD = \nabla \cdot \vB = 0$
identically; for even $(l, m, n)$ it is also periodic on the unit
cube with period $1$, so the same reference serves both periodic
and fully-physical boundary configurations in the study below.
The waveguide mode has nontrivial structure in every direction —
perfect for exercising all parts of the scheme.

\subsection{Using the analytic as SAT boundary data}

When \texttt{[source] type = "te\_waveguide\_mode"}, the code routes the
analytic solution as the boundary data $g$ on \emph{every} physical
face (not just the lower-$z$ face as in the plane-wave and beam
modes). For each SAT-face point, the scheme evaluates the analytic
at the local $(x, y, z, t)$ and passes it into the relevant
\texttt{APPLY\_SAT\_*} macro. The numerical state $U$ is then driven
toward the analytic through the incoming-characteristic projection
of every physical face; the outgoing characteristics are untouched,
so the test is consistent with the method's usual energy argument.

Concretely, the scheme being tested is:

\begin{itemize}
  \item Deep interior: 4th-order centered SBP stencil
    \texttt{D4CEN} in all three directions.
  \item Near-boundary rows on physical axes: SBP 4-2 one-sided
    closures (rows 0--3 at lower face, rows $N-4 \ldots N-1$ at
    upper face; see Sec.~\ref{sec:sbp-operators}).
  \item Exactly one boundary row per physical face: SAT penalty
    with $g$ evaluated from the analytic at that point.
  \item At edges (two physical faces meeting): two SAT penalties
    fire additively.
  \item At corners (three faces meeting): three SAT penalties
    fire additively.
  \item Time integration: classical explicit RK4.
\end{itemize}

If any of these pieces is implemented incorrectly, the convergence
rate at the affected configuration will fall below the theoretical
expectation. The decomposition below localises which class of
boundary structure is broken.

\subsection{The four configurations}

The boundary periodicities are varied per axis to change which
kinds of boundary structures are exercised:

\begin{center}
\begin{tabular}{lccccc}
\toprule
Configuration     & \texttt{periodic\_x} & \texttt{periodic\_y} & \texttt{periodic\_z}
                  & Physical faces / edges / corners & Expected $L^2$ rate \\
\midrule
\texttt{fully\_periodic} & true  & true  & true  & $0 / 0 / 0$   & $4$ \\
\texttt{z\_physical}     & true  & true  & false & $2 / 0 / 0$   & $3$ \\
\texttt{yz\_physical}    & true  & false & false & $4 / 4 / 0$   & $3$ \\
\texttt{all\_physical}   & false & false & false & $6 / 12 / 8$  & $3$ \\
\bottomrule
\end{tabular}
\end{center}

\paragraph{What each configuration localises.}
\begin{itemize}
  \item \texttt{fully\_periodic}: no SBP closures, no SAT penalties.
    All derivatives are \texttt{D4CEN} with periodic ghost zones.
    Tests only the interior stencil. Expected rate: $4$.
  \item \texttt{z\_physical}: only the two $z$-axis faces are
    physical; $x, y$ stay periodic.  Tests the interior plus the two
    $z$-face SBP closures plus the SAT penalty at the single
    lower-$z$ and upper-$z$ boundary rows.  No physical edges or
    corners are present. A drop in the rate here indicates a problem
    with either the SBP closure or the SAT penalty on a face.
  \item \texttt{yz\_physical}: adds four $y$-$z$ edges (where the
    physical $y$ and $z$ faces meet). A rate drop moving from
    \texttt{z\_physical} to \texttt{yz\_physical} indicates the
    additive edge penalty is mis-handled.
  \item \texttt{all\_physical}: the full setup — six faces, twelve
    edges, eight corners. A rate drop at this step would indicate
    the three-face corner penalty is mis-handled.
\end{itemize}

\subsection{Why the expected rates are $4$ and $3$}

The scheme is formally:
\begin{itemize}
  \item 4th-order in the interior (the \texttt{D4CEN} stencil is
    $O(h^4)$).
  \item 2nd-order locally at SBP-4-2 boundary rows (the Strand
    closures are $O(h^2)$ at each of the first/last four rows).
\end{itemize}
For a hyperbolic system discretized by an SBP operator with interior
order $2p$ and boundary truncation of order $s$, the classical
Svärd-Nordström result predicts a global $L^2$ convergence rate of
$\min(2p, s + 1)$ at late times; for SBP-4-2 with $2p = 4$ and
$s = 2$ this gives rate $3$.  The interior-only configuration
(periodic everywhere) sidesteps the boundary closures and retains
the full $4$th-order rate.

Two practical notes on these predictions:

\begin{itemize}
  \item The $s + 1$ rule applies to each physical face. Edges and
    corners do not further degrade the rate: the additive SAT
    penalty means each boundary contributes a separate energy
    estimate, and the rates combine as a maximum of errors, not
    a product.  This is the feature being tested.
  \item Measured rates often approach $3$ from above at finite $N$
    — rates of $3.1$--$3.2$ at $N = 48$ are typical for the SBP-4-2
    scheme; they asymptote to exactly $3$ as $N \to \infty$.  A rate
    persistently \emph{below} $3$ at moderate $N$ indicates a real
    problem.
\end{itemize}

\subsection{Running the convergence test}

The driver \texttt{scripts/convergence\_test.py} generates a TOML
file per $(\text{configuration}, \text{resolution})$ pair, runs the
solver, parses \texttt{l2\_norm.dat} for the final-time error, and
fits the convergence rate between consecutive resolutions.

\begin{lstlisting}[language=bash]
python scripts/convergence_test.py \
    --solver src/build/maxwell_system \
    --resolutions 16 24 32 48 \
    --final-time 0.1 \
    --np 4
\end{lstlisting}

Options:
\begin{itemize}
  \item \texttt{--solver PATH}: path to the built \texttt{maxwell\_system}
    executable.
  \item \texttt{--resolutions N1 N2 ...}: list of per-axis grid sizes
    to run. Default: \texttt{16 24 32 48}. A cubic grid $N \times N
    \times N$ is used for every run.
  \item \texttt{--final-time T}: physical simulation time per run
    (default $0.1$). Large enough that a few wave periods elapse;
    small enough that the error is dominated by the discretization
    rather than accumulated round-off.
  \item \texttt{--cfl C}: CFL factor (default $0.2$).
  \item \texttt{--np N}: number of MPI ranks (default $2$; $4$ is
    recommended for the $N = 48$ run to keep wall-clock under a
    minute).
  \item \texttt{--work-dir DIR}: scratch directory for per-run
    TOMLs and HDF5 outputs.
\end{itemize}

The driver runs the four configurations in sequence.  Each
configuration's section prints the per-resolution errors and the
pairwise rates:

\begin{verbatim}
=== all_physical  (expected L2 rate 3.0) ===
  N= 16 ... L2 = 0.01092
  N= 24 ... L2 = 0.003261
  N= 32 ... L2 = 0.001328
  N= 48 ... L2 = 0.000368
  resolution                16          24          32          48
  L2 error             0.01092    0.003261    0.001328    0.000368
  rate N_k,N_k+1          None       2.981       3.122       3.166
\end{verbatim}

A summary table compares the measured final rates against their
expected values:

\begin{verbatim}
=== SUMMARY ===
config              N=16        N=24        N=32        N=48  final rate
fully_periodic  0.001041   0.0002078   6.597e-05   1.306e-05    3.994   OK
z_physical       0.00739    0.002273   0.0009391    0.000261    3.158   OK
yz_physical     0.009489    0.002867    0.001174    0.000325    3.167   OK
all_physical     0.01092    0.003261    0.001328    0.000368    3.166   OK
\end{verbatim}

A rate is flagged as \texttt{SUSPICIOUS} if it deviates from the
expected value by more than $1.0$, and the driver exits nonzero.
The current code has all four rates at their expected values.

\subsection{Baseline results (as of the validated build)}

The numbers above are the validated baseline on the code in its
current state. For future regression checks:

\begin{center}
\begin{tabular}{lcccc|c}
\toprule
Configuration       & $N=16$   & $N=24$    & $N=32$    & $N=48$    & Final rate \\
\midrule
\texttt{fully\_periodic} & $1.04\cdot10^{-3}$ & $2.08\cdot10^{-4}$ &
  $6.60\cdot10^{-5}$ & $1.31\cdot10^{-5}$ & $3.99$ \\
\texttt{z\_physical} & $7.39\cdot10^{-3}$ & $2.27\cdot10^{-3}$ &
  $9.39\cdot10^{-4}$ & $2.61\cdot10^{-4}$ & $3.16$ \\
\texttt{yz\_physical} & $9.49\cdot10^{-3}$ & $2.87\cdot10^{-3}$ &
  $1.17\cdot10^{-3}$ & $3.25\cdot10^{-4}$ & $3.17$ \\
\texttt{all\_physical} & $1.09\cdot10^{-2}$ & $3.26\cdot10^{-3}$ &
  $1.33\cdot10^{-3}$ & $3.68\cdot10^{-4}$ & $3.17$ \\
\bottomrule
\end{tabular}
\end{center}

Run parameters: $T_\text{final} = 0.1$, $\mathrm{cfl} = 0.2$,
$\tau = 1.0$, four MPI ranks, uniform $\varepsilon = \mu = 1$,
$\sigma = 0$, analytic mode $(l, m, n) = (2, 2, 2)$.

\subsection{Interpreting a failed test}

If a future change to the boundary code causes one of the rates to
drop, the configuration hierarchy localises the fault:

\begin{itemize}
  \item \texttt{fully\_periodic} rate drops below $4$: the interior
    \texttt{D4CEN} stencil, the ghost-zone sync, or the periodic
    wrap is broken.
  \item \texttt{z\_physical} rate drops below $3$ while
    \texttt{fully\_periodic} stays at $4$: the $z$-face SBP
    closures or the lower-/upper-$z$ SAT penalty is broken.
  \item \texttt{yz\_physical} rate drops below $3$ while
    \texttt{z\_physical} is fine: the additive edge penalty at the
    $y$-$z$ edges is broken. (Since the $y$ and $z$ face SATs are
    each known to work in isolation from the previous row, the edge
    additivity is the new thing being tested.)
  \item \texttt{all\_physical} rate drops below $3$ while
    \texttt{yz\_physical} is fine: the three-face corner penalty at
    $x$-$y$-$z$ corners is broken.
\end{itemize}

This localisation is exactly the purpose of the test matrix. A
successful pass on all four rows confirms that every part of the
algorithm — interior stencil, face closure, face SAT, edge additivity,
corner additivity — is correctly implemented.

\paragraph{A subtle point: why the oblique-reflection failure mode
of Sec.~\ref{sec:oblique-reflection} does not break this test.}
The \texttt{te\_waveguide\_mode} has non-trivial structure in every
direction, and hits the $\pm x, \pm y$ faces at oblique angles — the
very configuration that is mis-handled by the $g = 0$ SAT used in the
lens demo. One could reasonably expect spurious reflections from the
$x$ and $y$ faces to pollute the measured rate.

This does not happen because the convergence test uses
$g = \text{analytic}$ on every physical face, not $g = 0$. The
absorbing-boundary failure mode only arises when the prescribed
incoming data disagrees with the true physical solution at the
boundary: with $g = 0$ and a nonzero field tangent to the wall, the
SAT pulls the wave toward zero instead of toward its true value, and
the resulting mismatch is the spurious reflection.  When $g$ equals
the exact analytic at every boundary point, the SAT imposes the
\emph{correct} incoming data on every face, regardless of the
incidence angle.  Mathematically, the SAT penalty is a weak-form
Dirichlet imposition of the incoming characteristic;
with the correct right-hand side it has no angular blind spot.

The convergence test therefore verifies that the SBP-SAT scheme,
\emph{given the correct incoming data}, converges to the analytic at
the expected rate. It is a test of the scheme's implementation, not
of the quality of absorbing boundaries in the lens-demo sense.

\paragraph{Caveats.}
\begin{itemize}
  \item The \texttt{te\_waveguide\_mode} source requires
    $\varepsilon = \mu = 1$ and $\sigma = 0$ (set in
    \texttt{[material.background]} by the driver). The test does not
    exercise the spatially-varying permittivity code path; that is a
    separate physics question, not a convergence question.
  \item The constraint fields $\Psi_D$, $\Psi_B$ are identically
    zero in the analytic. The constraint-damping blocks of the SAT
    are exercised only by their coupling to the EM blocks through
    residual numerical constraint violations, so a bug localised
    entirely to the $\Psi_D$, $\Psi_B$ blocks might not be caught
    at full sensitivity by this test. Production deployment with
    non-vacuum materials or with sources is the right check for
    those blocks.
  \item The test uses $g = \text{analytic}$ on every face and
    therefore does \emph{not} measure the quality of the $g = 0$
    absorbing boundary on oblique waves — that is a physics
    question about absorbing boundaries (Sec.~\ref{sec:oblique-reflection})
    and a design question about replacing the SAT with a PML
    (Sec.~\ref{sec:pml}), not a convergence question about the SAT
    implementation.
\end{itemize}

\section{References and notes}

\begin{itemize}
\item Strand, B. (1994), \emph{Summation by parts for finite difference
  approximations for $d/dx$}, J. Comput. Phys. 110, 47--67.
\item Svärd, M., Nordström, J. (2014), \emph{Review of summation-by-parts
  schemes for initial-boundary-value problems}, J. Comput. Phys. 268,
  17--38. The global-$L^2$ rate estimate
  $\min(2p, s+1)$ used in Sec.~\ref{sec:convergence} is from this
  review.
\item Munz, C.-D., Omnes, P., Schneider, R., Sonnendrücker, E., Voss, U.
  (2000), \emph{Divergence correction techniques for Maxwell solvers
  based on a hyperbolic model}, J. Comput. Phys. 161, 484--511.
\item B\'erenger, J.-P. (1994), \emph{A perfectly matched layer for the
  absorption of electromagnetic waves}, J. Comput. Phys. 114,
  185--200. The original split-field PML paper.
\item Roden, J.~A., Gedney, S.~D. (2000), \emph{Convolutional PML
  (CPML): an efficient FDTD implementation of the CFS-PML for
  arbitrary media}, Microwave Opt. Technol. Lett. 27, 334--339.
  The CPML variant recommended for the proposed extension in
  Sec.~\ref{sec:pml}.
\end{itemize}

\appendix

\section{Summary of SAT formulas}
\label{app:sat-formulas}

For reference, the full set of SAT contributions applied at the
boundary rows, derived in Sec.~\ref{sec:sat}. Let
$s = \tau \cdot (H^{-1}_0)/h$ be the per-axis penalty scale
(with $H^{-1}_0 = 48/17$ for the $(4,2)$ SBP norm), and define $c$,
$PP$ as in Sec.~\ref{sec:characteristics}. Where $g$ is zero, the
$\delta$ quantities reduce to the evolved fields themselves.

\paragraph{Lower-$z$ (nonzero $g$, the incoming beam):}
\begin{align*}
  \dot B_x   &\mathrel{-}= s \cdot \tfrac{c}{2} (\delta B_x - PP \cdot \delta D_y) \\
  \dot D_y   &\mathrel{-}= s \cdot \tfrac{c}{2} (-\delta B_x / PP + \delta D_y) \\
  \dot B_y   &\mathrel{-}= s \cdot \tfrac{c}{2} (\delta B_y + PP \cdot \delta D_x) \\
  \dot D_x   &\mathrel{-}= s \cdot \tfrac{c}{2} (\delta B_y / PP + \delta D_x) \\
  \dot B_z   &\mathrel{-}= s \cdot \tfrac{1}{2} (\delta B_z - \delta \Psi_B) \\
  \dot \Psi_B &\mathrel{-}= s \cdot \tfrac{1}{2} (-\delta B_z + \delta \Psi_B) \\
  \dot D_z   &\mathrel{-}= s \cdot \tfrac{1}{2} (\delta D_z - \delta \Psi_D) \\
  \dot \Psi_D &\mathrel{-}= s \cdot \tfrac{1}{2} (-\delta D_z + \delta \Psi_D)
\end{align*}

\paragraph{Upper-$z$ ($g = 0$):}
\begin{align*}
  \dot B_x   &\mathrel{-}= s \cdot \tfrac{c}{2} (B_x + PP \cdot D_y) \\
  \dot D_y   &\mathrel{-}= s \cdot \tfrac{c}{2} (B_x / PP + D_y) \\
  \dot B_y   &\mathrel{-}= s \cdot \tfrac{c}{2} (B_y - PP \cdot D_x) \\
  \dot D_x   &\mathrel{-}= s \cdot \tfrac{c}{2} (-B_y / PP + D_x) \\
  \dot B_z   &\mathrel{-}= s \cdot \tfrac{1}{2} (B_z + \Psi_B) \\
  \dot \Psi_B &\mathrel{-}= s \cdot \tfrac{1}{2} (B_z + \Psi_B) \\
  \dot D_z   &\mathrel{-}= s \cdot \tfrac{1}{2} (D_z + \Psi_D) \\
  \dot \Psi_D &\mathrel{-}= s \cdot \tfrac{1}{2} (D_z + \Psi_D)
\end{align*}

\paragraph{Lower-$x$, Lower-$y$, Upper-$x$, Upper-$y$:} Obtained from the
lower-$z$ and upper-$z$ formulas above by cyclic permutation of the axes.
The explicit forms are in \texttt{maxwell\_eqs.c}, macros
\texttt{APPLY\_SAT\_\{LOWER,UPPER\}\_\{X,Y\}}.

\end{document}
